\documentclass[aspectratio=169]{beamer}
\usetheme{metropolis}
\metroset{sectionpage=none, numbering=counter, progressbar=frametitle}
\usepackage{amsmath, amssymb}
\usepackage[T1]{fontenc}

\title{Domain-Partition Algorithm for Turbomachinery Blades}
\subtitle{Slide 10: Transfinite Interpolation --- Coons Patch}
\date{\today}
\author{}
\institute{}

\begin{document}

\begin{frame}{Transfinite Interpolation --- Coons Patch}

\begin{center}
\Large
\begin{aligned}
X(u,v) ={}& (1-v)\cdot S(u) + v\cdot N(u) \\
&+ (1-u)\cdot W(v) + u\cdot E(v) \\
&- \bigl[(1-u)(1-v)\cdot c_{00} + u(1-v)\cdot c_{10} \\
&\quad + (1-u)v\cdot c_{01} + uv\cdot c_{11}\bigr]
\end{aligned}
\end{center}

\vspace{0.2em}
\begin{itemize}
    \item \textbf{Bilinear interpolation} from 4 boundary curves
    \item \textbf{Guarantees $C^0$ continuity} at block edges
    \item \textbf{Cell counts solved conforming:} opposite sides match
\end{itemize}

\vspace{0.2em}
\small
\textit{Code:} \texttt{tmesh\_partition.py:\_coons (lines 881--892)} \quad
\textit{Lit:} Coons 1967, ``Surfaces for Computer-Aided Design''

\note{
Speaker notes:
Transfinite interpolation, also known as the Coons patch, constructs a surface from four boundary curves. S and N are the south and north boundaries, W and E are the west and east boundaries. The first four terms perform linear blending along each boundary, while the subtracted bilinear corner term removes the double-counted corner contributions. This guarantees C0 continuity across adjacent blocks because shared boundaries are interpolated identically. In our implementation, the _coons function in tmesh_partition.py at lines 881 through 892 computes this on a regular u-v grid. The cell counts on opposite sides of each block are solved to be conforming, meaning the number of cells along a shared edge matches on both sides. The method was introduced by Steven Coons in 1967 in his report ``Surfaces for Computer-Aided Design of Space Forms.''
}

\end{frame}

\end{document}
